\documentclass[11pt,a4paper]{article}

\usepackage[margin=1in]{geometry}
\usepackage{amsmath,amssymb,amsfonts,amsthm}
\usepackage{booktabs}
\usepackage{array}
\usepackage{hyperref}
\usepackage{xcolor}

\hypersetup{colorlinks=true,linkcolor=blue,citecolor=blue,urlcolor=blue}

\newcommand{\CP}{\mathbb{CP}}
\newcommand{\as}{\alpha_s}
\newcommand{\Nf}{N_f}
\newcommand{\Nc}{N_c}

\newtheorem{theorem}{Theorem}
\newtheorem{lemma}[theorem]{Lemma}
\newtheorem{proposition}[theorem]{Proposition}
\theoremstyle{definition}
\newtheorem{definition}[theorem]{Definition}
\theoremstyle{remark}
\newtheorem{remark}[theorem]{Remark}

\definecolor{pass}{RGB}{0,120,0}
\definecolor{fail}{RGB}{180,0,0}
\definecolor{warn}{RGB}{180,120,0}

\begin{document}

\title{Rigorous Verification of $A_b = 1/42$:\\
Does $1/(\Nf \times b_0)$ Follow from QCD Running?}

\author{DFD Research Audit}

\date{23 March 2026}

\maketitle

\begin{abstract}
We rigorously verify the claim in the DFD v67 mass dictionary that
$A_b = 1/42 = 1/(\Nf \times b_0)$ with $\Nf = 6$ and
$b_0 = (11\Nc - 2\Nf)/3 = 7$.  The arithmetic factorization is confirmed.
However, numerical computation of the actual 1-loop QCD mass running from
$M_P$ to $m_b$ yields a ratio of $\approx 3.8$--$3.9$ (i.e., the running
mass \emph{increases} by a factor of $\sim 4$ from Planck to $m_b$ scale),
not a suppression of $1/42 \approx 0.024$.  The value $A_b = 1/42$ is the
\emph{measured} $b$-quark Yukawa coupling $m_b/(v/\sqrt{2}) = 4180/174100
= 1/41.65 \approx 1/42$, expressed via QCD parameters through pattern
identification.  It is not derived from the QCD renormalization group.
Similarly, $A_s = 6/7 = \Nf/b_0$ is an empirical fit accurate to $0.03\%$,
not a running ratio.  Both are \textbf{pattern-level} inputs, not
theorem-grade derivations.
\end{abstract}

\tableofcontents

%==========================================================================
\section{The Claim Under Examination}
\label{sec:claim}
%==========================================================================

The DFD v67 charged-fermion mass dictionary uses the master formula
\begin{equation}\label{eq:master}
m_f = A_f \,\alpha^{n_f}\, \frac{v}{\sqrt{2}}\,,
\end{equation}
with the down-type quark prefactors organized via a ``QCD running operator''
\begin{equation}\label{eq:Qd}
Q_d = \mathrm{diag}\!\left(1,\;\frac{6}{7},\;\frac{1}{42}\right)
\quad\text{on}\quad \mathcal{H}_{\mathrm{gen}}\,.
\end{equation}

The claimed derivation is:
\begin{itemize}
\item $b_0 = (11\Nc - 2\Nf)/3$ is the 1-loop QCD $\beta$-function coefficient.
\item With $\Nc = 3$, $\Nf = 6$: $b_0 = (33 - 12)/3 = 21/3 = 7$.
\item The $b$-quark entry: $Q_d[3,3] = 1/(\Nf \times b_0) = 1/(6 \times 7) = 1/42$.
\item The $s$-quark entry: $Q_d[2,2] = \Nf/b_0 = 6/7$.
\item The $d$-quark entry: $Q_d[1,1] = 1$ (no running modification).
\end{itemize}

\noindent\textbf{Questions to resolve:}
\begin{enumerate}
\item Is the arithmetic $b_0 = 7$ correct? (\S\ref{sec:arithmetic})
\item Is $A_b = 1/42$ derivable from 1-loop QCD mass running? (\S\ref{sec:running})
\item What does the numerical QCD running actually give? (\S\ref{sec:numerical})
\item What is the actual origin of $1/42$? (\S\ref{sec:origin})
\item Same analysis for $A_s = 6/7$. (\S\ref{sec:strange})
\item Is $\alpha_s$ at the relevant scales predicted by DFD or taken as external input? (\S\ref{sec:external})
\end{enumerate}

%==========================================================================
\section{Arithmetic Verification of $b_0 = 7$}
\label{sec:arithmetic}
%==========================================================================

\textcolor{pass}{\textbf{PASS.}}

The 1-loop QCD $\beta$-function coefficient in the standard convention is
\begin{equation}
b_0 = \frac{11\Nc - 2\Nf}{3}\,.
\end{equation}

With $\Nc = 3$ (QCD color gauge group $SU(3)$) and $\Nf = 6$ (all six
quark flavors active):
\begin{equation}
b_0 = \frac{11 \times 3 - 2 \times 6}{3} = \frac{33 - 12}{3} = \frac{21}{3} = 7\,.
\end{equation}

Therefore $\Nf \times b_0 = 6 \times 7 = 42$, confirming
\begin{equation}
\frac{1}{\Nf \times b_0} = \frac{1}{42}\,.
\end{equation}

\begin{remark}[Convention note]
Some references define $\beta_0 = (11\Nc - 2\Nf)/(12\pi)$ or
$\beta_0 = (11\Nc - 2\Nf)$, absorbing factors of $4\pi$ differently.
The DFD convention $b_0 = (11\Nc - 2\Nf)/3$ matches the most common
textbook convention (Peskin--Schroeder, PDG).  With this convention,
the 1-loop running coupling is
\[
\as(\mu) = \frac{\as(\mu_0)}{1 + \frac{b_0}{2\pi}\,\as(\mu_0)\,\ln(\mu/\mu_0)}\,.
\]
\end{remark}

\begin{remark}[$\Nf = 6$ assumption]
The choice $\Nf = 6$ means all six quarks are treated as active.
In standard QCD, $\Nf = 6$ applies only above the top threshold
$\mu > m_t \approx 173$~GeV.  At the $b$-quark scale ($\mu \sim 4$~GeV),
only $\Nf = 5$ flavors are active ($b_0 = 23/3$).  This is a
\textbf{physically questionable} choice if $b_0 = 7$ is meant to describe
running at the $b$ scale, but numerologically convenient.
\end{remark}

%==========================================================================
\section{The Standard QCD Mass Running Formula}
\label{sec:running}
%==========================================================================

The 1-loop running of quark masses in QCD (in the $\overline{\mathrm{MS}}$ scheme) is
\begin{equation}\label{eq:mass-running}
\frac{m(\mu_1)}{m(\mu_2)} = \left[\frac{\as(\mu_1)}{\as(\mu_2)}\right]^{d_0}\,,
\qquad
d_0 = \frac{\gamma_0}{2b_0}\,,
\end{equation}
where:
\begin{itemize}
\item $\gamma_0 = 8$ is the 1-loop quark mass anomalous dimension (universal in QCD),
\item $b_0 = (11\Nc - 2\Nf)/3$ is the 1-loop $\beta$-function coefficient.
\end{itemize}

For $\Nf = 6$: $d_0 = 8/(2 \times 7) = 4/7 \approx 0.571$.

For $\Nf = 5$: $d_0 = 8/(2 \times 23/3) = 12/23 \approx 0.522$.

\subsection{What the running ratio would need to be}

For $A_b = 1/42$ to be a running ratio, we would need
\begin{equation}
\left[\frac{\as(m_b)}{\as(M_P)}\right]^{4/7} = \frac{1}{42} \approx 0.0238\,.
\end{equation}

This requires
\begin{equation}
\frac{\as(m_b)}{\as(M_P)} = \left(\frac{1}{42}\right)^{7/4} = 42^{-7/4} \approx 1.4 \times 10^{-3}\,,
\end{equation}
i.e., $\as(M_P)/\as(m_b) \sim 700$.  Since $\as(m_b) \approx 0.21$ and
$\as(M_P) \approx 0.019$, the actual ratio is $\as(m_b)/\as(M_P) \approx 11$,
nowhere near the required $\sim 700$.

\textcolor{fail}{\textbf{FAIL: The QCD running cannot produce $1/42$.}}

%==========================================================================
\section{Numerical Computation of QCD Mass Running}
\label{sec:numerical}
%==========================================================================

\subsection{Running $\as$ from $M_Z$ to $M_P$ and $m_b$}

Starting from $\as(M_Z) = 0.1179$ (PDG 2024), using 1-loop running with
flavor thresholds:

\begin{center}
\begin{tabular}{lccc}
\toprule
Scale $\mu$ & $\Nf$ & $\as(\mu)$ & Method \\
\midrule
$m_b = 4.18$~GeV & 5 & 0.2118 & Run down from $M_Z$ \\
$M_Z = 91.19$~GeV & 5 & 0.1179 & Input (PDG) \\
$m_t = 172.8$~GeV & 5 & 0.1080 & Run up from $M_Z$ \\
$M_P = 1.22 \times 10^{19}$~GeV & 6 & 0.0191 & Run up from $m_t$ \\
\bottomrule
\end{tabular}
\end{center}

\subsection{Mass running ratio $m(m_b)/m(M_P)$}

\subsubsection{Naive: $\Nf = 6$ throughout (as DFD implicitly assumes)}

\begin{equation}
\frac{m(m_b)}{m(M_P)} = \left[\frac{\as(m_b)}{\as(M_P)}\right]^{4/7}
= \left[\frac{0.2118}{0.0192}\right]^{4/7} \approx \mathbf{3.94}\,.
\end{equation}

\subsubsection{Proper: with flavor threshold at $m_t$}

\begin{equation}
\frac{m(m_b)}{m(M_P)} =
\underbrace{\left[\frac{\as(m_b)}{\as(m_t)}\right]^{12/23}}_{\approx\,1.42}
\times
\underbrace{\left[\frac{\as(m_t)}{\as(M_P)}\right]^{4/7}}_{\approx\,2.69}
\approx \mathbf{3.83}\,.
\end{equation}

\subsection{Summary of running results}

\begin{center}
\begin{tabular}{lcc}
\toprule
Quantity & Value & Compare to $1/42$ \\
\midrule
$m(m_b)/m(M_P)$ (naive, $\Nf=6$) & 3.94 & $\times\, 166$ too large \\
$m(m_b)/m(M_P)$ (proper, thresholds) & 3.83 & $\times\, 161$ too large \\
$1/42$ & 0.0238 & --- \\
\bottomrule
\end{tabular}
\end{center}

\noindent The QCD mass running from $M_P$ to $m_b$ \emph{increases} the mass
by a factor of $\sim 4$ (due to the growing $\as$), not a suppression by 42.
The running and the claimed value differ by two orders of magnitude.

\textcolor{fail}{\textbf{VERDICT: $A_b = 1/42$ is NOT the QCD mass running
ratio from $M_P$ to $m_b$.}}

%==========================================================================
\section{The Actual Origin of $1/42$}
\label{sec:origin}
%==========================================================================

\subsection{$1/42$ is the measured Yukawa coupling}

With $n_b = 0$, the mass formula reduces to
\begin{equation}
m_b = A_b \times \frac{v}{\sqrt{2}} = A_b \times 174.1\;\text{GeV}\,.
\end{equation}

Solving for $A_b$:
\begin{equation}
A_b = \frac{m_b}{v/\sqrt{2}} = \frac{4.180}{174.1} = 0.02401 = \frac{1}{41.65} \approx \frac{1}{42}\,.
\end{equation}

The value $1/42$ is simply the \textbf{measured $b$-quark Yukawa coupling}
(up to the $\sqrt{2}$ convention), rounded to the nearest integer denominator.
The error is
\begin{equation}
\frac{1/42 - 1/41.65}{1/41.65} = \frac{41.65 - 42}{42} = -0.83\%\,.
\end{equation}

\subsection{The factorization $42 = \Nf \times b_0$ is numerological}

The decomposition $42 = 6 \times 7$ where $6 = \Nf$ and
$7 = (11\Nc - 2\Nf)/3$ is algebraically valid.  However:
\begin{enumerate}
\item No RG equation produces $1/(\Nf \cdot b_0)$ as a mass ratio or
  Yukawa coupling at any scale.
\item The standard QCD mass running gives a completely different value
  ($\sim 3.8$, not $1/42$).
\item The source files (see \texttt{Af\_Derivation\_Trace.tex},
  \S\S\,8--9) explicitly classify this as ``pattern-level, not
  theorem-grade.''
\item The \texttt{Af\_RIGOROUS\_COMPUTATION.py} audit script marks $Q_d$
  entries as ``pattern identifications, not derived from RG flow.''
\end{enumerate}

\subsection{Could there be an indirect path through DFD-specific RG?}

One might hypothesize that DFD modifies the running (e.g., via the
density field coupling), producing a non-standard RG flow.  However:
\begin{itemize}
\item No such modified RG has been derived in any DFD document.
\item The v67 dictionary and Appendix~K explicitly cite standard QCD
  parameters $\Nf$ and $b_0$, not DFD-modified ones.
\item If the running were DFD-modified, $b_0 = 7$ (the standard QCD value)
  would not be the correct coefficient.
\end{itemize}

%==========================================================================
\section{Verification of $A_s = 6/7 = \Nf/b_0$}
\label{sec:strange}
%==========================================================================

The strange quark has $n_s = 3/2$, so
\begin{equation}
A_s = \frac{m_s}{\alpha^{3/2}\,(v/\sqrt{2})} = \frac{93\;\text{MeV}}{(1/137.036)^{3/2} \times 174100\;\text{MeV}}
= \frac{93}{108.53} = 0.8569\,.
\end{equation}

Comparing to $\Nf/b_0 = 6/7 = 0.8571$: match to $0.03\%$.

As with $A_b$, this is the \textbf{measured mass ratio} expressed via QCD
parameters.  The QCD running from $M_P$ to $m_s$ does not produce $6/7$;
it gives a factor of $\sim 4$--5 (similar to the $b$ quark, with additional
threshold effects from $c$ and $b$).

\subsection{The $Q_d$ pattern structure}

The three entries form a suggestive but unexplained pattern:
\begin{align}
Q_d[1,1] &= 1 \quad (\text{identity, no modification})\,, \\
Q_d[2,2] &= \frac{\Nf}{b_0} = \frac{6}{7}\,, \\
Q_d[3,3] &= \frac{1}{\Nf \cdot b_0} = \frac{1}{42}\,.
\end{align}

The ratio between consecutive entries is
\begin{equation}
\frac{Q_d[2,2]}{Q_d[1,1]} = \frac{6}{7}\,,\qquad
\frac{Q_d[3,3]}{Q_d[2,2]} = \frac{1/42}{6/7} = \frac{1}{36} = \frac{1}{\Nf^2}\,.
\end{equation}

This means the jump from gen~2 to gen~3 is $1/\Nf^2 = 1/36$, while from
gen~1 to gen~2 it is $\Nf/b_0 = 6/7$.  These are structurally different
ratios, arguing against a uniform running mechanism.

%==========================================================================
\section{Is $\as$ at the Relevant Scales Predicted by DFD?}
\label{sec:external}
%==========================================================================

\textcolor{warn}{\textbf{CRITICAL POINT:}} The DFD framework derives the
\emph{fine-structure constant} $\alpha = 1/137.036$ from the Chern--Simons
level sum with $k_{\max} = 60$.  However, the QCD coupling $\as$ is a
\textbf{separate coupling constant} with its own RG flow.

\begin{itemize}
\item DFD does \textbf{not} derive $\as$ at any scale from microsector
  geometry.
\item DFD does \textbf{not} derive the QCD $\beta$-function; $b_0 = 7$ is
  standard perturbative QCD.
\item The mass dictionary uses $\alpha$ (QED), not $\as$ (QCD), in the
  exponent $\alpha^{n_f}$.
\end{itemize}

If $A_b = 1/42$ were genuinely derived from QCD running, it would require
$\as$ at two scales as input.  Since DFD does not predict $\as$, this would
constitute \textbf{external SM input}, not a prediction from the microsector.

However, as shown in \S\ref{sec:numerical}, $1/42$ does not come from QCD
running at all, so this question is moot: $Q_d$ is an empirical operator,
not a derived one.

%==========================================================================
\section{Comparison: Model A vs Model B for the $b$ Quark}
\label{sec:models}
%==========================================================================

Two parametrizations exist in the DFD literature:

\begin{center}
\begin{tabular}{lcccc}
\toprule
Model & $n_b$ & $A_b$ & $m_b^{(\text{pred})}$ & Error \\
\midrule
v67 (Model A) & 0 & $1/42$ & $(1/42) \times 174.1 = 4145$~MeV & $-0.83\%$ \\
Bundle (Model B) & 1 & $\pi$ & $\pi \times (1/137.036) \times 174100 = 3991$~MeV & $-4.51\%$ \\
\bottomrule
\end{tabular}
\end{center}

Model~A is 5.4$\times$ more accurate but uses $A_b = 1/42$ as an empirical
input.  Model~B is geometrically derived ($\pi$ from $S^3$ color integration)
but less accurate.

The resolution proposed in \texttt{B\_Quark\_Exponent\_Derivation.tex} is
``color-connected vertex saturation'' (Path~4): 3rd-generation quarks have
$n_f \to n_f - 1$ due to color saturation, giving $n_b = 0$ while
$A_b = 1/42$ arises from the $S^3$ color integration at the saturated vertex.
This is physically motivated but not yet rigorously derived.

%==========================================================================
\section{Conclusions}
\label{sec:conclusions}
%==========================================================================

\begin{enumerate}

\item \textcolor{pass}{\textbf{Arithmetic: CONFIRMED.}}
$b_0 = (11 \times 3 - 2 \times 6)/3 = 7$ and $1/(\Nf \times b_0) = 1/42$.

\item \textcolor{fail}{\textbf{QCD running derivation: FAILS.}}
The 1-loop QCD mass running from $M_P$ to $m_b$ gives $m(m_b)/m(M_P) \approx 3.8$--$3.9$,
not $1/42 \approx 0.024$.  These differ by a factor of $\sim 160$.

\item \textcolor{warn}{\textbf{True origin: empirical Yukawa coupling.}}
$A_b = 1/42 \approx m_b/(v/\sqrt{2}) = 4180/174100$.  The factorization
$42 = \Nf \times b_0$ is a numerological coincidence involving QCD
parameters, not a derived result.

\item \textcolor{warn}{\textbf{Similarly: $A_s = 6/7$ is empirical.}}
$A_s = \Nf/b_0 = 6/7 \approx m_s/(\alpha^{3/2}\,v/\sqrt{2})$, matching to
$0.03\%$.  Again a pattern identification, not a running ratio.

\item \textcolor{warn}{\textbf{Classification: pattern-level input.}}
The $Q_d$ operator entries are among the \textbf{least rigorously derived}
components of the v67 dictionary.  They are empirical fits expressed in QCD
language.  The Af derivation trace audit classifies them as ``pattern-level,
not theorem-grade.''

\item \textbf{What would be needed for a genuine derivation:}
\begin{itemize}
\item An explicit RG computation (within DFD or standard QCD) showing that
  the down-type Yukawa coupling, when run from the Planck scale to the
  weak/quark scale with all threshold corrections, yields exactly $\Nf/b_0$
  for gen~2 and $1/(\Nf \cdot b_0)$ for gen~3.
\item Alternatively: a microsector operator derivation showing that $Q_d$
  emerges from $A_5$ or $\CP^2$ geometry with entries involving $b_0$.
\item Neither exists in the current DFD corpus.
\end{itemize}

\end{enumerate}

\subsection*{Bottom line for the mass dictionary}

$A_b = 1/42$ produces an excellent prediction ($-0.83\%$ error) and has a
striking decomposition in QCD parameters.  But calling it a ``QCD running
factor'' is \textbf{misleading}.  It is an empirical input that \emph{happens}
to factorize as $1/(\Nf \times b_0)$.  Whether this coincidence has a deeper
origin is an open question.

\end{document}
